\section{Encryption}

Public key encryption schemes enable confidential communication by allowing
anyone possessing a public key to encrypt messages that can only be decrypted
using the corresponding secret key.

\paragraph{Interfaces.} A public key encryption scheme consists of the
following algorithms:

\begin{itemize}

  \item $(\pk, \sk) \sample \PKEgen(\secparam)$ : the key generation algorithm
    outputs public and secret key pair $(\pk, \sk)$.

  \item $\ciph \leftarrow \PKEencrypt(\pk, \msg)$ : the encryption algorithm
    takes as input a public key $\pk$ and a message $\msg$ and outputs
    ciphertext $\ciph$.

  \item $\msg \leftarrow \PKEdecrypt(\sk, \ciph)$ : the decryption algorithm
    takes as input a secret key $\sk$ and ciphertext $\ciph$ and outputs
    plaintext $\msg$.

\end{itemize}

\paragraph{Correctness.} An encryption scheme must satisfy the following
correctness requirement: for every key pair $(\pk, \sk)$ generated by
$\PKEgen(\secparam)$ and every message $\msg$, it holds that $\PKEdecrypt(\sk,
\PKEencrypt(\pk, \msg)) = \msg$.

\paragraph{IND-CCA security.} We require public key encryption schemes to
satisfy indistinguishability under chosen ciphertext attack (IND-CCA), defined
by the security experiment defined as follows:
%
\begin{enumerate} 

\item The challenger generates $(\pk, \sk) \sample
  \PKEgen(\secparam)$ and gives $\pk$ to $\adv$.

  \item The adversary is given oracle access to $\deco(\ciph) =
    \PKEdecrypt(\sk,\ciph)$.

  \item The adversary outputs two equal-length messages $(\msg_0,\msg_1)$.

  \item The challenger samples $b \sample \{0,1\}$ and computes the challenge
    ciphertext $\ciph^\star \sample \PKEencrypt(\pk,\msg_b)$. They send
    $\ciph^\star$ to the adversary.

  \item The adversary retains access to $\deco$, but may not query
    $\ciph^\star$. They eventually output a bit $b'$.

\end{enumerate}
%
The adversary wins if $b' = b$. Their IND-CCA advantage is defined as follows: 
\begin{equation*}
\Adv^{\mathsf{ind-cca}}_{\PKE, \ \adv}(\secparam) = \left| \ \Pr[b'=b] - \frac{1}{2} \ \right|.
\end{equation*}
%
The public key encryption scheme is IND-CCA secure if this advantage is small.